\section*{Capítulo \ref{ch:b1:lewis-resonance-vsepr} --- Estructuras de Lewis, resonancia y modelo RPECV}

\begin{solution}{exo:b1:lewis-resonance-vsepr:1}
\ce{H2O}: dos enlaces \ce{O-H} y dos \omterm{def:b1:lewis-resonance-vsepr:lewis}{pares no enlazantes} en O. \ce{NH3}:
tres enlaces \ce{N-H} y un \omterm{def:b1:lewis-resonance-vsepr:lewis}{par no enlazante} en N. \ce{CO2}: \ce{O=C=O},
dos \omterm{def:b1:lewis-resonance-vsepr:lewis}{pares no enlazantes} en cada O. \ce{HCN}: \ce{H-C#N}, un \omterm{def:b1:lewis-resonance-vsepr:lewis}{par no
enlazante} en N. \ce{CH2O}: carbono unido a dos H y con enlace doble con
O, dos \omterm{def:b1:lewis-resonance-vsepr:lewis}{pares no enlazantes} en O. En estas cinco, todas las \omterm{def:b1:lewis-resonance-vsepr:formal-charge}{cargas
formales} son nulas. \ce{NH4+}: cuatro enlaces \ce{N-H}, ningún \omterm{def:b1:lewis-resonance-vsepr:lewis}{par no
enlazante}; nitrógeno $5 - 0 - 4 = +1$, la carga del ion.
\end{solution}

\begin{solution}{exo:b1:lewis-resonance-vsepr:2}
\ce{H2S}: \ce{AX2E2}, angular, menos de $109.5^\circ$ (medido:
$92^\circ$). \ce{PCl3}: \ce{AX3E}, piramidal trigonal, unos $100^\circ$.
\ce{BF3}: \ce{AX3}, trigonal plana, $120^\circ$. \ce{SiH4}: \ce{AX4},
tetraédrica, $109.5^\circ$. \ce{CO3^{2-}}: \ce{AX3} (el enlace doble
cuenta una vez), trigonal plana, $120^\circ$.
% ledger: geo:H2S.aHSH, geo:PCl3.aClPCl
\end{solution}

\begin{solution}{exo:b1:lewis-resonance-vsepr:3}
Apolares: \ce{CCl4} (\ce{AX4} tetraédrica), \ce{BF3} (\ce{AX3}) y
\ce{CO2} (\ce{AX2} lineal): los \omterm{def:b1:lewis-resonance-vsepr:dipole}{dipolos de enlace} se anulan. Polares:
\ce{CH2Cl2} (dos tipos de enlace distintos en un tetraedro), \ce{NH3}
(piramidal) y \ce{SO2} (angular).
\end{solution}

\begin{solution}{exo:b1:lewis-resonance-vsepr:4}
\ce{H+} es el \omterm{def:b1:lewis-resonance-vsepr:lewis-acid}{ácido de Lewis}, y \ce{H2O}, la base: el nuevo enlace
\ce{O-H} de \ce{H3O+} es dativo en el momento de formarse (después, los
tres enlaces \ce{O-H} son idénticos). \ce{Cu^{2+}} es el ácido y cada
\ce{NH3}, una base: los cuatro enlaces \ce{Cu-N} son dativos.
\end{solution}

\begin{solution}{exo:b1:lewis-resonance-vsepr:5}
\chemfig{CH_3-C(=[:60]O)-[:-60]\chemabove{O}{\scriptstyle\ominus}}
$\leftrightarrow$
\chemfig{CH_3-C(-[:60]\chemabove{O}{\scriptstyle\ominus})=[:-60]O}.
Las dos estructuras son equivalentes: los dos enlaces \ce{C-O} tienen
orden 1.5 y la misma longitud, entre la del enlace doble y la del
sencillo. En el ácido metanoico los dos oxígenos son distintos (uno lleva
el H), las estructuras no son equivalentes y los enlaces conservan
longitudes distintas, 120.2 (\ce{C=O}) y \qty{134.3}{pm} (\ce{C-OH}).
\end{solution}

\begin{solution}{exo:b1:lewis-resonance-vsepr:6}
$N = 6 + 4 \times 6 + 2 = 32$ electrones. Con cuatro enlaces sencillos y
tres \omterm{def:b1:lewis-resonance-vsepr:lewis}{pares no enlazantes} en cada oxígeno: azufre $6 - 0 - 4 = +2$, cada
oxígeno $6 - 6 - 1 = -1$; total, $+2 - 4 = -2$. Forma \ce{AX4},
tetraédrica. Con dos enlaces \ce{S=O}, la elección de cuáles de los
cuatro oxígenos llevan el enlace doble da $\binom{4}{2} = 6$ estructuras.
\end{solution}

\begin{solution}{exo:b1:lewis-resonance-vsepr:7}
\ce{SF4}: \ce{AX4E}, de balancín. \ce{ClF3}: \ce{AX3E2}, en forma de T.
\ce{XeF2}: \ce{AX2E3}, lineal. \ce{XeF4}: \ce{AX4E2}, plano-cuadrada. En
una bipirámide trigonal, una posición axial tiene tres vecinos a
$90^\circ$, y una ecuatorial, solo dos: los voluminosos \omterm{def:b1:lewis-resonance-vsepr:lewis}{pares no
enlazantes} van donde sufren menos repulsiones a $90^\circ$.
\end{solution}

\begin{solution}{exo:b1:lewis-resonance-vsepr:8}
El azufre y el fósforo son más grandes y menos electronegativos que el
oxígeno y el nitrógeno: los \omterm{def:b1:lewis-resonance-vsepr:lewis}{pares enlazantes} están más lejos del átomo
central y atraídos hacia el hidrógeno, de modo que se repelen menos, y
los \omterm{def:b1:lewis-resonance-vsepr:lewis}{pares no enlazantes} empujan los enlaces hacia $90^\circ$.
\end{solution}

\begin{solution}{exo:b1:lewis-resonance-vsepr:9}
\ce{CH3-C#N}: el carbono del \ce{CH3} es sp$^3$; el carbono del nitrilo,
sp; el nitrógeno, sp; 5 enlaces $\sigma$ (tres \ce{C-H}, \ce{C-C} y uno en
\ce{C#N}) y 2 enlaces $\pi$. \ce{CH2=CH-CH3}: los dos carbonos del
alqueno son sp$^2$, y el carbono del metilo, sp$^3$; 8 enlaces $\sigma$
(seis \ce{C-H} y dos \ce{C-C}) y 1 enlace $\pi$.
\end{solution}

\begin{solution}{exo:b1:lewis-resonance-vsepr:10}
Las dos moléculas son piramidales, con un \omterm{def:b1:lewis-resonance-vsepr:lewis}{par no enlazante} en el
nitrógeno. El \omterm{def:b1:lewis-resonance-vsepr:lewis}{par no enlazante} da una contribución que apunta (de lo
negativo a lo positivo) del par hacia el núcleo. En \ce{NH3}, el
nitrógeno es el extremo negativo de cada enlace: los \omterm{def:b1:lewis-resonance-vsepr:dipole}{dipolos de enlace}
apuntan de N hacia los hidrógenos, en el mismo sentido que la
contribución del \omterm{def:b1:lewis-resonance-vsepr:lewis}{par no enlazante}, y se suman. En \ce{NF3}, el flúor es
el extremo negativo: los \omterm{def:b1:lewis-resonance-vsepr:dipole}{dipolos de enlace} apuntan de los flúores hacia
N, en contra de la contribución del \omterm{def:b1:lewis-resonance-vsepr:lewis}{par no enlazante}, y ambos casi se
anulan.
\end{solution}

\begin{solution}{exo:b1:lewis-resonance-vsepr:11}
$\mu_b = \mu/(2\cos(\theta/2)) = 1.63/(2\cos 59.75^\circ) =
\qty{1.62}{\debye} = \qty{5.40e-30}{C.m}$. Entonces, $\delta = \mu_b/(e\,d) =
5.40 \times 10^{-30}/(1.602 \times 10^{-19} \times 143.2 \times 10^{-12})
= 0.24$.
\end{solution}

\begin{solution}{exo:b1:lewis-resonance-vsepr:12}
\ce{N=N=O} con \omterm{def:b1:lewis-resonance-vsepr:formal-charge}{cargas formales} $-1$, $+1$, $0$; \ce{N#N-O} con $0$,
$+1$, $-1$. El enlace \ce{N-N} (\qty{112.8}{pm}) está cerca del enlace
triple de \ce{N2} (\qty{109.8}{pm}): la estructura \ce{N#N-O} pesa más, y
sitúa la carga negativa sobre el oxígeno, el átomo más electronegativo.
Sin embargo, el enlace \ce{N-O} es más corto que un enlace sencillo: la
otra estructura también contribuye.
\end{solution}

\begin{solution}{pb:b1:lewis-resonance-vsepr:1}
\textbf{1.} \ce{N2} 10, \ce{NO} 11, \ce{NO2} 17, \ce{N2O} 16, \ce{O3} 18,
\ce{HNO3} 24.
\textbf{2.} \ce{N#N} con un \omterm{def:b1:lewis-resonance-vsepr:lewis}{par no enlazante} en cada nitrógeno.
\textbf{3.} \ce{N=O} con dos \omterm{def:b1:lewis-resonance-vsepr:lewis}{pares no enlazantes} en el oxígeno, y un \omterm{def:b1:lewis-resonance-vsepr:lewis}{par
no enlazante} y un electrón solitario en el nitrógeno. Con 11 electrones,
un número impar, siempre hay un átomo con una cuenta impar: el nitrógeno
tiene 7 electrones a su alrededor.
\textbf{4.} \ce{N=N=O}: $-1$, $+1$, $0$; \ce{N#N-O}: $0$, $+1$, $-1$.
\textbf{5.} \ce{O=O-O}: oxígeno central $+1$ (un \omterm{def:b1:lewis-resonance-vsepr:lewis}{par no enlazante}, tres
enlaces), el oxígeno terminal con enlace sencillo $-1$, el otro 0.
\textbf{6.} Como en el \cref{ex:b1:lewis-resonance-vsepr:hno3}: nitrógeno
$+1$, el oxígeno con un enlace sencillo y sin hidrógeno $-1$, los demás 0.
\textbf{7.} \ce{NO} y \ce{NO2} (11 y 17 electrones).
\textbf{8.} Dos estructuras equivalentes, con el enlace doble a uno u
otro lado: orden 1.5 para cada enlace.
\textbf{9.} Sí: \qty{127.8}{pm} está entre el enlace doble
(\qty{120.8}{pm}) y el sencillo (\qty{147.5}{pm}).
\textbf{10.} El enlace doble, sucesivamente, en cada uno de los tres
oxígenos; cada enlace \ce{N-O} es doble en una estructura de cada tres:
orden $4/3$.
\textbf{11.} El enlace largo, \qty{140.6}{pm}, es \ce{N-OH}, un enlace
sencillo. Los dos oxígenos sin hidrógeno comparten un enlace $\pi$ a
través de dos \omterm{def:b1:lewis-resonance-vsepr:resonance}{formas resonantes} equivalentes: dos enlaces de orden 1.5,
de 121.1 y \qty{119.9}{pm}.
\textbf{12.} El nitrógeno lleva un electrón solitario en lugar de un \omterm{def:b1:lewis-resonance-vsepr:lewis}{par
no enlazante}: repele los \omterm{def:b1:lewis-resonance-vsepr:lewis}{pares enlazantes} menos de lo que lo haría un
par, y los enlaces se abren por encima de $120^\circ$.
\textbf{13.} \ce{N#N-O}, con la carga negativa en el oxígeno.
\textbf{14.} \ce{O3}: \ce{AX2E}, angular. \ce{N2O}: \ce{AX2} (N central),
lineal. \ce{NO3-}: \ce{AX3}, trigonal plana. Nitrógeno de \ce{HNO3}:
\ce{AX3}, trigonal plana.
\textbf{15.} El \omterm{def:b1:lewis-resonance-vsepr:lewis}{par no enlazante} del oxígeno central repele los \omterm{def:b1:lewis-resonance-vsepr:lewis}{pares
enlazantes} más de lo que estos se repelen entre sí.
\textbf{16.} El ozono es angular y su átomo central difiere de los
extremos (\omterm{def:b1:lewis-resonance-vsepr:formal-charge}{carga formal} positiva en el centro, carga negativa repartida
entre los extremos): un pequeño dipolo a lo largo de la bisectriz.
\ce{N2O} es lineal, pero sus dos extremos son átomos distintos: un
pequeño dipolo. \ce{N2} es simétrica: ninguno.
\textbf{17.} \ce{CO2} es lineal (\ce{AX2}): sus \omterm{def:b1:lewis-resonance-vsepr:dipole}{dipolos de enlace} se
anulan. \ce{SO2} es angular (\ce{AX2E}): se suman.
\textbf{18.} \ce{AX3E}: el \omterm{def:b1:lewis-resonance-vsepr:lewis}{par no enlazante} cierra los ángulos por debajo
de $109.5^\circ$; medido: $106.7^\circ$.
\textbf{19.} \ce{AX2E2}: dos \omterm{def:b1:lewis-resonance-vsepr:lewis}{pares no enlazantes} lo cierran aún más;
$104.5^\circ$.
\textbf{20.} $\mu = 2\mu_{OH}\cos(\theta/2)$
(\cref{prop:b1:lewis-resonance-vsepr:dipole-sum}).
\textbf{21.} $\mu_{OH} = 1.857/(2\cos 52.24^\circ) = 1.857/1.2246 =
\qty{1.52}{\debye}$.
\textbf{22.} $1.516 \times 3.336 \times 10^{-30} = \qty{5.06e-30}{C.m}$.
\textbf{23.} $e\,d = 1.602 \times 10^{-19} \times 95.8 \times 10^{-12} =
\qty{1.535e-29}{C.m}$, es decir, \qty{4.60}{\debye}.
\textbf{24.} $\delta = 5.06 \times 10^{-30}/1.535 \times 10^{-29} =
\textbf{0.33}$: el enlace \ce{O-H} del agua lleva un tercio de la carga
que llevaría un enlace totalmente iónico.
\end{solution}
